% DRAFT (2026-11-28) for the coordinator's read. Not published.
\documentclass[11pt]{article}
\usepackage[margin=1.1in]{geometry}
\usepackage{amsmath,graphicx,xcolor}
\usepackage{hyperref}
\hypersetup{colorlinks=true,urlcolor=blue!50!black,linkcolor=blue!50!black}
\setlength{\parskip}{0.6em}\setlength{\parindent}{0pt}
\begin{document}
\begin{center}
{\Large\bfseries Chaos around a ``bumpy'' black hole: a plain summary}\\[0.4em]
Sumit Pal (AI-assisted) \quad$\cdot$\quad 28 November 2026 (draft)
\end{center}

\textbf{The question.} Einstein's theory predicts that a spinning black hole is described by one exact solution, the
Kerr metric. Orbits around a Kerr black hole are remarkably orderly: besides energy and angular momentum there is a
third, hidden conserved quantity (Carter's constant), so every orbit is predictable for all time. Physicists test
whether real black holes are Kerr by asking what would change if the black hole were slightly ``bumpy'', with an
extra bulge or flattening of its mass. The Manko--Novikov spacetime is an exact solution of Einstein's equations
that does exactly this: a number $\beta$ measures the bump, and $\beta=0$ is ordinary Kerr.

\textbf{What we showed.} For a bumpy black hole of this kind, the orbits of freely falling particles near a
particular unstable orbit are \emph{chaotic} in the precise mathematical sense: there is a ``horseshoe''---a tangle
in which orbits can follow any prescribed sequence of two kinds of loops---and the system has positive
\emph{topological entropy}, a standard measure of how fast orbits become unpredictable. We prove this
\begin{itemize}
\item for a specific bump $\beta=3/100$, by a \textbf{computer-assisted proof}, and
\item for \emph{every} sufficiently small bump of either sign, by a pencil-and-paper argument (the Melnikov method)
whose one key number was computed rigorously.
\end{itemize}
The same computer code, run on the ordinary Kerr black hole, correctly finds \emph{no} chaos there.

\textbf{What ``computer-assisted proof'' means here.} The computer never trusts a rounded number. Every quantity is
carried as an interval guaranteed to contain the true value (``interval arithmetic''), so the final inequalities are
mathematically certain, provided the software that does the interval arithmetic is correct. To guard against errors,
the key decisions were recomputed in a second, independent arithmetic library; the equations of motion were
re-derived from the published formulas by a separate program that had not seen ours; and every decision rule and
expected outcome was written down before each computation.

\textbf{Why it matters.} Chaos had been seen in this spacetime in computer simulations before, but simulations cannot
prove chaos: rounding errors can create or hide it. To our knowledge this is the first computer-assisted proof of chaos
for orbits in the Manko--Novikov spacetime. It supports the picture that any deviation from Kerr of this kind breaks the
orderly motion, which matters for gravitational-wave tests of black holes.

\textbf{What we do \emph{not} claim.} We give no number for how chaotic the motion is (no entropy value). We treat one
black hole (mass 5, spin 3 in our units), one family of orbits and one energy. For the ``every small bump'' result we
do not know how small ``small'' is. The orbit calculations rely on one validated-integration library (CAPD); a second
one has not been run.

\textbf{How this was done.} The work was carried out with substantial help from AI systems (Claude by Anthropic and
Gemini by Google), under the author's direction. All the ``independent'' checks above were done by separate AI sessions
or systems, \emph{not} by human experts; no external expert has yet read the proofs. The full paper, the code, and a
guide for checking every claim are in the record; the software is archived at
\href{https://doi.org/10.5281/zenodo.24602464}{doi:10.5281/zenodo.24602464}.
\end{document}
